\chapter{Background and related work}
\label{chap:background}

This chapter introduces the concepts on which FlowSync is built. Section~\ref{sec:scheduling}
defines the scheduling problems considered and the objectives used to compare plans, and
Section~\ref{sec:graph} presents the graph representation that underlies both the exact methods
and the learned ones. Sections~\ref{sec:cpsat} and~\ref{sec:dispatching} describe the two
classical families of solution methods, constraint programming and dispatching rules, which in
FlowSync play the role of teacher and of baseline respectively. Section~\ref{sec:ml-scheduling}
reviews the use of machine learning for scheduling and Section~\ref{sec:llm} the use of language
models to build the problem description. Section~\ref{sec:positioning} summarises how these
pieces are combined in FlowSync.

\status{This chapter reviews the literature and does not depend on the progress of the project.
Where it mentions the role of a method in FlowSync, the last section and
Table~\ref{tab:positioning} say whether that role is already implemented.}

\section{Scheduling problems}
\label{sec:scheduling}

Scheduling is the allocation of limited resources to activities over time. The field has a
standard classification, due to Graham and described in detail by Pinedo~\cite{pinedo2022}, that
describes a problem by the machine environment (how many resources and how the activities flow
through them), the constraints and the objective. The environments relevant to this work are
the following:

\begin{itemize}
  \item in the \textbf{job shop}, each job is a chain of operations, each operation must be
        executed on one specific machine, and a machine processes one operation at a time;
  \item in the \textbf{flexible job shop}, each operation can be executed by any machine of a
        given set, so that the problem includes also the choice of the machine
        (\emph{routing})~\cite{brandimarte1993};
  \item in the \textbf{resource-constrained project scheduling problem} (RCPSP), the activities
        are linked by general precedence relations, instead of chains, and consume units of
        resources that can be shared among several activities at the same time.
\end{itemize}

On top of these basic problems, real applications add \emph{sequence-dependent setup times} (the
time to prepare a resource depends on the pair of consecutive activities), working
\emph{calendars} and \emph{unavailabilities}, deadlines and priorities. These are the features
that FlowSync includes in its common model (Section~\ref{sec:common-model}).

\subsection{Formal model}

Let $A$ be the set of activities and $R$ the set of resources. Each activity $a \in A$ belongs to
a work order $j(a)$ and has a set of eligible resources $R_a \subseteq R$; if it is executed on
$r \in R_a$, it lasts $p_{a,r}$. Let $P \subseteq A \times A$ be the set of precedences. A
\emph{schedule} assigns to each activity a resource $\rho(a) \in R_a$ and a start time $s_a$ such
that:
\begin{align}
  s_b &\ge s_a + p_{a,\rho(a)} && \text{for each } (a,b) \in P, \label{eq:prec}\\
  s_b &\ge s_a + p_{a,r} + \sigma_r(a,b) \quad\text{or}\quad
  s_a \ge s_b + p_{b,r} + \sigma_r(b,a), \label{eq:disj}
\end{align}
where~\eqref{eq:disj} holds for each pair of distinct activities $a, b$ with $\rho(a) = \rho(b) = r$,
and $\sigma_r(a,b)$ is the setup time of resource $r$ between $a$ and $b$ (zero if absent).
In addition, an activity can be processed only while its resource is available according to its
calendar. Constraint~\eqref{eq:prec} states that an activity cannot start before its
predecessors are finished; constraint~\eqref{eq:disj} states that two activities on the same
resource cannot overlap, and that setups must fit between them. In the RCPSP the second
constraint is replaced by a capacity limit on the sum of the resource units used at each instant.

\subsection{Objectives}

The quality of a schedule is measured by one or more objective functions. Let $C_j$ be the
completion time of work order $j$, $d_j$ its due date and $w_j$ its weight. The objectives
considered in this thesis are:
\begin{itemize}
  \item the \textbf{makespan} $C_{\max} = \max_j C_j$, the time at which the last activity ends;
  \item the \textbf{total weighted tardiness} $\sum_j w_j T_j$, where $T_j = \max(0, C_j - d_j)$;
  \item the \textbf{overtime}, the amount of work executed outside the regular shifts;
  \item the \textbf{total setup time}.
\end{itemize}
In FlowSync these objectives are combined with weights that depend on the domain profile and that
the planner can change, so that the same engine can serve different priorities.

\subsection{Complexity and benchmarks}

Scheduling problems are in general hard. The job shop is NP-hard even in simple
cases~\cite{garey1976}, and the RCPSP is NP-hard in the strong sense~\cite{blazewicz1983}, so no
algorithm is known that is guaranteed to find the optimal schedule in a time that grows
polynomially with the size of the instance. For this reason, the research community
compares methods on shared collections of instances. The benchmarks most relevant for this work
are Taillard's instances for flow shop, job shop and open shop problems~\cite{taillard1993}, the
ten flexible job shop instances of Brandimarte~\cite{brandimarte1993}, and the PSPLIB library for
the RCPSP~\cite{kolisch1997}, which contains sets of projects of 30, 60, 90 and 120 activities.
Most of these benchmarks target the makespan and do not include setups or calendars. For this
reason FlowSync uses them for validation and testing, while the instances used for training are
generated synthetically, with the features of the common model.

\section{Disjunctive graph representation}
\label{sec:graph}

A scheduling instance can be represented as a graph, the \emph{disjunctive graph}, introduced for
the job shop by Balas~\cite{balas1969}. The nodes are the activities, plus a source node $S$ and
a sink node $T$. There are two kinds of arcs:
\begin{itemize}
  \item \textbf{conjunctive arcs} represent the precedences, and are always present; they connect
        $S$ to the first activities and the last activities to $T$;
  \item \textbf{disjunctive arcs} connect each pair of activities that compete for the same
        resource. They are not oriented: they express that the two activities must be executed
        one before the other, but not which one comes first.
\end{itemize}
Figure~\ref{fig:disjunctive} shows a small example with two jobs and three machines.

\begin{figure}[t]
  \centering
  \resizebox{\linewidth}{!}{%
  \begin{tikzpicture}[
      act/.style={circle, draw, minimum size=8mm, inner sep=0pt, font=\small},
      term/.style={circle, draw, fill=black!8, minimum size=8mm, inner sep=0pt, font=\small},
      conj/.style={-{Stealth[length=2mm]}, thick},
      disjA/.style={{Stealth[length=2mm]}-{Stealth[length=2mm]}, dashed},
      disjB/.style={{Stealth[length=2mm]}-{Stealth[length=2mm]}, densely dotted}]
    \node[term] (S) at (-1.6,0.8) {$S$};
    \node[act]  (a1) at (0,1.6) {$a_1$};
    \node[act]  (a2) at (3,1.6) {$a_2$};
    \node[act]  (a3) at (6,1.6) {$a_3$};
    \node[act]  (b1) at (0,0) {$b_1$};
    \node[act]  (b2) at (3,0) {$b_2$};
    \node[term] (T) at (7.6,0.8) {$T$};
    \draw[conj] (S) -- (a1);  \draw[conj] (a1) -- (a2); \draw[conj] (a2) -- (a3); \draw[conj] (a3) -- (T);
    \draw[conj] (S) -- (b1);  \draw[conj] (b1) -- (b2); \draw[conj] (b2) -- (T);
    \draw[disjA] (a1) -- node[pos=0.28, above, sloped, font=\scriptsize] {$M_1$} (b2);
    \draw[disjB] (a2) -- node[pos=0.28, below, sloped, font=\scriptsize] {$M_2$} (b1);
  \end{tikzpicture}}
  \caption{Disjunctive graph of a job shop with two jobs. Solid arrows are the precedences
    within the jobs ($a_1 \to a_2 \to a_3$ and $b_1 \to b_2$). The dashed and dotted double arrows
    link the pairs of activities that share a machine: $a_1$ and $b_2$ on $M_1$, $a_2$ and $b_1$
    on $M_2$; $a_3$ is alone on $M_3$ and needs no disjunctive arc. Choosing an orientation for each
    pair produces a schedule.}
  \label{fig:disjunctive}
\end{figure}

To build a schedule, one has to choose an orientation for every disjunctive arc, a
\emph{selection}. A selection is feasible if the resulting graph has no cycle, and in that case
the earliest start time of each activity is given by the longest path from $S$ to the activity;
the longest path from $S$ to $T$ is the makespan. The scheduling problem becomes the search for
the selection that minimises the objective. The representation also extends to the flexible job
shop, where each operation is in addition linked to the machines that can execute it, and to
setups, which add a weight to the arcs.

This representation matters for two reasons. First, it makes explicit that scheduling is a
\emph{graph} problem, so that its instances can be processed by graph-based models. Second, it
describes the schedule as a sequence of decisions, one arc at a time, which is the view adopted
by the learned dispatching policies discussed in Section~\ref{sec:ml-scheduling}. In FlowSync the
same structure is built step by step by the dispatcher of the AI module (Section~\ref{sec:m2}).

\section{Exact methods: constraint programming}
\label{sec:cpsat}

\emph{Constraint programming} (CP) solves a problem by describing it with variables and
constraints and by searching systematically for the solutions, using propagation to prune the
search space. Scheduling is one of its most successful applications, thanks to dedicated
modelling elements: an \emph{interval variable} represents an activity with its start, duration
and end; a \emph{no-overlap} constraint states that the intervals on a resource are disjoint, and
a \emph{cumulative} constraint limits the use of shared resources. Constraints~\eqref{eq:prec}
and~\eqref{eq:disj} map directly onto these elements.

The solver used in this work is CP-SAT, part of Google's OR-Tools suite~\cite{ortools}. It
combines constraint propagation with techniques from Boolean satisfiability and with linear
programming relaxations~\cite{perron2023cpsat}, and it is widely used as a strong baseline for
scheduling. Given a time limit, CP-SAT returns a status:
\texttt{OPTIMAL} when it has proved that the solution found is the best one,
\texttt{FEASIBLE} when it has a valid solution but could not prove optimality,
\texttt{INFEASIBLE} when it has proved that no solution exists, and \texttt{UNKNOWN} when it has
found nothing within the limit.

Exact solvers have clear advantages: they handle many kinds of constraints in a uniform way and
they can prove optimality. Their limit is the computation time, which grows quickly with the
size of the instance, and which is not predictable. For this reason FlowSync is designed not to
call CP-SAT when it serves a request. The solver is used offline, to compute good schedules on
synthetic instances: these schedules are the examples from which the neural model will learn
(Section~\ref{sec:ml-scheduling}). In this role CP-SAT is the \emph{teacher}.

\section{Dispatching rules}
\label{sec:dispatching}

A \emph{dispatching rule} builds the schedule step by step. Whenever a resource becomes free, the
rule computes a priority index for each activity that is ready (all its predecessors are done and
it is eligible for the resource) and selects the one with the best index. The schedule is built
in a single pass, and by construction it respects precedences and capacities. The most common
rules are:
\begin{itemize}
  \item \textbf{EDD} (earliest due date): selects the activity of the order with the nearest
        deadline;
  \item \textbf{SPT} (shortest processing time): selects the shortest activity, which tends to
        reduce the number of waiting orders;
  \item \textbf{ATC} (apparent tardiness cost): introduced by Vepsäläinen and
        Morton~\cite{vepsalainen1987} for job shops with weighted tardiness. In its basic form,
        the index of activity $a$ at time $t$ is
        \[
          I_a(t) \;=\; \frac{w_a}{p_a}\,\exp\!\left(-\,\frac{\max(d_a - p_a - t,\,0)}{K\,\bar p}\right),
        \]
        where $w_a$, $p_a$ and $d_a$ are the weight, the duration and the due date associated
        with the activity, $\bar p$ is the average duration of the waiting activities and $K$ a
        look-ahead parameter. The index favours short and important activities, and it grows
        quickly when the due date approaches.
\end{itemize}
Dispatching rules are instantaneous, simple to implement and to explain, and they scale to very
large instances. Their weakness is that each decision depends only on the current state and on a
fixed formula, so that the quality of the schedule can be far from the optimum, and different
rules are best on different instances. In FlowSync, EDD and ATC have two roles: they are the
\emph{baselines} that the learned model has to beat, and they are designed to be the
\emph{fallback} used when the model cannot provide a plan.

\section{Machine learning for scheduling}
\label{sec:ml-scheduling}

Machine learning has been applied to combinatorial optimisation in several ways. The survey of
Bengio, Lodi and Prouvost~\cite{bengio2021} distinguishes between approaches that use a model to
make decisions directly, those that use it to support a classical algorithm (for example by
suggesting which branch to explore), and those that use it to configure the algorithm. The
learned dispatching policies discussed here belong to the first group: a neural network replaces
the hand-written priority formula of the previous section.

\subsection{Learned dispatching policies}

The idea is to keep the structure of a dispatcher, and to let a neural network compute the
priorities. Zhang et al.~\cite{zhang2020l2d} represent the state of a job shop as a disjunctive
graph, process it with a graph neural network and train, by reinforcement learning, a policy that
chooses the next operation to dispatch. Since the network works on the graph, it does not depend on
the size of the instance and, as reported by the authors, it generalises to instances larger than
those seen in training. Song et al.~\cite{song2023fjsp} extend the approach to the flexible job
shop: the state is a heterogeneous graph with operation and machine nodes, and the policy
chooses jointly the operation and the machine. Graph neural networks~\cite{scarselli2009}
compute, for each node, a vector that summarises the node and its neighbourhood by exchanging
messages along the arcs: this is the tool that makes the disjunctive graph usable as input of a
neural model.

\subsection{Imitation learning}

The approaches above learn by reinforcement learning, in which the model explores and receives a
reward at the end of the schedule. A different path is \emph{imitation learning}, in which the
model learns to reproduce the decisions of an expert~\cite{hussein2017imitation}. In its simplest
form, \emph{behaviour cloning}, the expert decisions are collected as pairs (state, decision) and
the model is trained with supervised learning to predict the decision from the state. This is
simpler and more stable than reinforcement learning, because the signal is dense: every decision
has a label. Its main limitation is that the model is exposed only to the states visited by the
expert, and at run time, after a mistake, it can reach states that it has never seen. Methods that
query the expert again on the states visited by the learner, such as DAgger~\cite{ross2011dagger},
address this problem.

FlowSync follows the imitation approach, with a specific choice of the expert: CP-SAT solves
synthetic instances offline, and the decisions that lead to the best schedule found are the
labels. The model will be limited by the quality of the teacher, but in exchange it needs no
reward design, and it can run at inference time without any solver.

\subsection{Estimation of durations and setups}

A second use of learning concerns the data. Durations and setup times are not known with
certainty, and planning with a single average value hides the risk of delays. They can be
estimated from the history of the executions with a regression model; in \emph{quantile
regression}~\cite{koenker1978} the model predicts, instead of the mean, a chosen quantile of the
distribution, so that an interval of values can be obtained and the confidence of the
estimate is explicit. Event logs of real processes, such as the public logs of the Dutch
municipalities~\cite{bpic2015} or of a traffic fine management process~\cite{roadfines}, are a source
of examples of durations with this kind of variability, and are the planned testbed for this
estimator in FlowSync.

\section{Language models for data entry}
\label{sec:llm}

Describing a scheduling problem requires many pieces of information (activities, durations,
precedences, resources, deadlines), and entering them through forms is slow and error-prone. Large
language models (LLMs) can read a request written in natural language and extract these
elements. The task is related to the \emph{auto-formulation} of optimisation models from text,
studied for instance by Ramamonjison et al.~\cite{ramamonjison2022} for linear programs: the
text is translated into a structured representation of the problem, which is then passed to a
solver.

Language models can however produce plausible but wrong output, and for this reason the design of
FlowSync gives the LLM a limited role. It does not decide the plan: it only fills a
\emph{problem specification} (ProblemSpec), whose structure is fixed by a schema. The specification is then
validated by deterministic code, which checks types, references and consistency (for example,
that every precedence refers to existing activities), and the planner is asked for the information
that is missing instead of guessing it. In this way an error of the language model
can produce an incomplete specification, which is detected, but not an infeasible schedule.

\section{Positioning of FlowSync}
\label{sec:positioning}

The previous sections describe the building blocks, and FlowSync combines them in a specific
way. Compared with the learned dispatching approaches of the literature, which target the job
shop or the flexible job shop with the makespan as objective, FlowSync works on a common model
that includes setups, calendars, deadlines and a weighted combination of objectives, and that is
meant to cover factory, public administration and project scheduling (the instance generator
covers the first two so far). Compared with exact methods, it moves the cost of optimisation to
an offline training phase, and keeps a classical dispatcher as fallback. Compared with a fully
automatic system, it places the planner at the end of the process, with three candidate plans and
their indicators. Table~\ref{tab:positioning} summarises the role of each method and how far it
is implemented.

\begin{table}[t]
  \centering
  \begin{tabular}{@{}p{3.1cm}p{5.2cm}p{4.2cm}@{}}
    \toprule
    Method & Role in FlowSync & State on \statusdate \\
    \midrule
    CP-SAT & teacher: produces the labelled schedules, offline & implemented \\
    EDD, ATC & baselines; fallback when the model fails & implemented as baselines \\
    Neural dispatching policy & optimisation engine, online & dataset ready, model to train \\
    Quantile regression & estimate of durations and setups & to do \\
    LLM & translation of the request into a ProblemSpec, validated & to do \\
    \bottomrule
  \end{tabular}
  \caption{Methods described in this chapter and their role in FlowSync.}
  \label{tab:positioning}
\end{table}
